\title{Direct Preference Optimization: Your Language Model is Secretly a Reward Model}

\begin{abstract}
Although large-scale unsupervised language models (LMs) acquire extensive knowledge about the world and demonstrate some capacity for reasoning, exerting fine-grained control over their outputs remains challenging, primarily due to the fully unsupervised nature of their training regimes. Current approaches to enhance steerability typically involve assembling human-annotated data that expresses relative preferences among model outputs, followed by fine-tuning the original unsupervised LM to reflect these human judgments—most commonly using reinforcement learning from human feedback (RLHF). Nevertheless, RLHF introduces considerable complexity and instability, as it requires constructing a reward model that captures human preferences and subsequently using reinforcement learning to adapt the LM so as to maximize the predicted reward, all while attempting to prevent significant deviation from the pretrained model. In this work, we propose a novel way to parameterize the reward model within RLHF, which facilitates the direct derivation of the optimal policy in closed form. This allows us to address the classical RLHF optimization problem using only a straightforward classification objective. The algorithm we introduce, termed \textit{Direct Preference Optimization} (DPO), offers robust performance and stability while being computationally efficient. DPO obviates the need for generation-time sampling from the LM during fine-tuning and reduces the necessity for extensive hyperparameter searches. Empirically, we find that DPO is capable of aligning language models with human preferences at least as effectively as, and sometimes better than, previously established techniques. Significantly, DPO surpasses PPO-based RLHF in regulating generation sentiment and delivers equal or improved results in tasks like summarization and single-turn dialogue, all while being notably easier to implement and train.
\end{abstract}

\section{Introduction}
Large unsupervised language models (LMs), when exposed to extensive datasets, demonstrate a range of remarkable abilities~\citep{chowdhery2022palm, brown2020language, touvron2023llama,bubeck2023sparks}. Nonetheless, such models are trained on data produced by humans who possess diverse intentions, preferences, and expertise. This heterogeneity means that not all behaviors present in the training corpus are suitable for emulation; for instance, it is useful for an AI coding assistant to \textit{recognize} frequent programming errors to correct them, but when producing code, it is preferable to orient the model towards the exemplary, albeit potentially uncommon, coding practices embedded in its data. Analogously, while a language model should be \textit{familiar} with misconceptions endorsed by a significant fraction of people, we do not wish the model to propagate these misconceptions in proportion to their prevalence in human belief. Consequently, distinguishing and eliciting \emph{desired responses and behaviors} from the broad spectrum of a model’s \textit{knowledge and capacities} is essential to ensure that AI systems are safe, effective, and aligned~\citep{ouyang2022training}. Whereas most contemporary approaches attempt to align LMs with human values by employing reinforcement learning (RL) techniques, we argue that the RL-based learning objective used in these settings can be optimized precisely through a straightforward binary cross-entropy loss, thereby streamlining the overall process of preference learning.

In broad terms, prevalent techniques shape the behavior of language models by leveraging carefully curated datasets that reflect human preferences regarding safe and beneficial conduct. This preference learning phase takes place subsequent to a massive, unsupervised pre-training stage over text corpora. While supervised fine-tuning on human-annotated demonstrations of exemplary responses represents the simplest form of preference learning, the most effective methods in practice have involved reinforcement learning from human (or AI) feedback (RLHF/RLAIF; \citep{christiano2017deep,bai2022constitutional}). RLHF approaches entail training a reward model on a collection of human preference judgments, and subsequently using RL to adapt the model policy such that it produces outputs evaluated as high-reward, all while constraining the policy to remain reasonably close to the original pre-trained model. Though RLHF has enabled the development of language models with strong conversational and coding competencies, it introduces substantial complexity beyond supervised techniques, necessitating the training of multiple models and iterative policy sampling, which collectively impose notable computational overhead.

In this work, we present a method for training language models to align with human preferences without relying on explicit reward models or reinforcement learning techniques. Our approach, called \textit{{Direct Preference Optimization} (DPO)}, offers an implicit means of optimizing the same underlying objective as typical RLHF methods—maximizing reward while constraining the KL-divergence from the initial policy—yet it is both easier to implement and more straightforward to train. Conceptually, DPO works by increasing the log odds of generating responses that are favored over those that are not; it does so by introducing an adaptive, per-example importance weighting mechanism, which we find mitigates the model collapse observed with naive probability ratio formulations. Similar to prior methods, DPO leverages a theoretical framework for modeling preferences, such as the Bradley-Terry model (\cite{bradley1952rankanalysis}), which assesses the agreement between a reward function and empirical preference judgements. In contrast to conventional pipelines that employ this model to derive a loss for training a separate reward network (followed by optimizing the main policy to maximize this reward), DPO applies a change of variables to directly relate the preference-based loss to the policy’s parameters. With access to datasets of human preferences on model outputs, DPO enables the optimization of the policy using a basic binary cross-entropy loss, yielding a policy that is optimal for a reward implicitly specified by the preferences.

The core contribution of this paper is Direct Preference Optimization (DPO), a straightforward algorithm for preference-based language model training that does not require reinforcement learning. Experimental results demonstrate that DPO matches or surpasses the performance of established approaches—including those based on PPO for RLHF—across various tasks such as sentiment control, text summarization, and conversational response generation, even for models as large as 6B parameters.

\section{Related Work}

As self-supervised language models have grown larger, they have demonstrated the capability to perform certain tasks in a zero-shot setting \citep{radford2019language} or by leveraging a small number of prompt examples \citep{gpt3,megatron,chowdhery2022palm}. Nonetheless, their effectiveness on specific applications and their alignment with user expectations can be markedly enhanced through fine-tuning on collections of instructional data accompanied by human-generated outputs \citep{mishra-etal-2022-cross,sanh2022multitask,chung2022scaling,thoppilan2022lamda}. This process, referred to as 'instruction-tuning,' allows large language models to extend their capabilities to novel instructions not seen during training, thereby boosting their overall usefulness \citep{chung2022scaling}. 

Although instruction-tuning has achieved notable progress, gathering comparative human assessments of output quality is often more feasible than sourcing expert-level demonstrations. As a result, subsequent research has focused on further fine-tuning LLMs using datasets built from human preference judgments, resulting in improvements in areas such as machine translation \citep{kreutzer-etal-2018-reliability}, text summarization \citep{stiennon2022learning,ziegler2020finetuning}, narrative generation \citep{ziegler2020finetuning}, and instruction adherence \citep{ouyang2022training,ramamurthy2023is}. The standard approach in these works involves initially training a neural reward model that aligns with observed preferences, frequently employing a preference modeling framework like the Bradley-Terry model \citep{bradley1952rankanalysis}. Subsequently, the language model is further optimized via reinforcement learning methods—such as REINFORCE \citep{williams1992reinforce}, proximal policy optimization (PPO; \cite{schulman2017proximal}), or their adaptations \citep{ramamurthy2023is}—to maximize the expected reward. 

An adjacent research direction utilizes LLMs, which have already been adapted for instruction following through human feedback, to synthesize additional preference data targeting specific qualities such as safety or non-maleficence \citep{bai2022constitutional}. Here, weak human supervision, often in the form of textual guidelines, steers the annotation process carried out by the LLM. Such strategies represent the intersection of two significant research streams: training language models with reinforcement learning to satisfy a variety of objectives~\citep{Ranzato2015SequenceLT,paulus2018a,wu2018learning}, and general methodologies for preference-based learning from human feedback \citep{christiano2017deep,kupcsik2018learning}. Despite the attractiveness of preference-based optimization, using reinforcement learning to fine-tune large language models poses substantial practical difficulties. In this context, the present work introduces a principled alternative for optimizing models according to relative preferences without relying on reinforcement learning.

Research on learning policies from preferences outside the realm of language has encompassed both bandit and reinforcement learning paradigms, with a variety of methods developed in the literature. One such setting, contextual dueling bandit (CDB; \cite{yue2012karmed,dudik2015contextual}), extends contextual bandit learning by using action preferences or rankings instead of scalar rewards. In scenarios lacking explicit reward signals, the theoretical framework for CDBs adopts the concept of a \textit{von Neumann winner}: a policy that has a minimum expected win probability of 50\% when compared with any alternative policy \citep{dudik2015contextual}. Notably, CDBs typically operate in an online regime where preference information arrives sequentially, in contrast to human preference learning, which often relies on a static, offline collection of action pairs annotated with preferences \citep{yan2022human}. In a related direction, \textit{preference-based RL} (PbRL) is concerned with learning from binary preferences derived from an opaque ‘scoring’ function, instead of observable rewards \citep{BusaFekete2014,ruiz2023dueling}. Numerous PbRL algorithms have been devised—including strategies that utilize off-policy preference data—but these frequently proceed by first inferring an implicit reward or scoring model and then optimizing it in a separate stage \citep{jain2013learning,BusaFekete2014,christiano2017deep,sadigh2017active,kupcsik2018learning}. In contrast, our approach proposes to learn a policy directly in one stage by optimizing it according to observed preferences.

\section{Preliminaries}\label{section:prelims}

We summarize the RLHF methodology as articulated in \citeauthor{ziegler2020finetuning} and later adopted in works such as \citep{stiennon2022learning, bai2022training, ouyang2022training}. The RLHF workflow generally comprises three principal stages: 1) supervised fine-tuning (SFT); 2) preference-based reward modeling; and 3) reinforcement learning optimization.

\textbf{SFT}: The process is initiated by adapting a pre-trained language model to specific downstream applications (such as dialogue systems or summarization) via supervised learning on curated, high-quality datasets, yielding a fine-tuned model $\pi^\text{SFT}$.

\textbf{Reward Modelling Phase}: Subsequently, the SFT model receives prompt inputs $x$ and generates response pairs $(y_1, y_2)\sim \pi^\text{SFT}(y \mid x)$. Human annotators are then tasked with indicating a preferred completion, denoted by $y_w\succ y_l \mid x$, where $y_w$ is favored over $y_l$ among the sampled responses. These human choices are modeled as being governed by an unobserved reward function $r^*(y, x)$. To characterize the preference data, models such as the Bradley-Terry (BT) model \cite{bradley1952rankanalysis} are commonly used; more general ranking approaches like the Plackett-Luce model \citep{plackett1975analysis, luce2012individual} can also be accommodated, especially when multiple ranked responses are available. Under the BT formulation, the distribution over human preferences $p^*$ takes the following form:

\begin{equation}\label{eq:bradley-terry}
    p^*(y_1\succ y_2 \mid x)=\frac{\exp\left(r^*(x, y_1)\right)}{\exp\left(r^*(x, y_1)\right) + \exp\left(r^*(x, y_2)\right)}.
\end{equation}

Given access to a fixed dataset of comparison data $\mathcal{D} = \bigl\{x^{(i)}, y_w^{(i)}, y_l^{(i)}\bigr\}_{i=1}^N$ drawn according to $p^*$, one can parameterize the reward function as $r_\phi(x, y)$ and learn its parameters through maximum likelihood estimation. Treating this as a binary classification problem, the associated loss is the negative log-likelihood:

\begin{equation}\label{eq:reward_model}
    \mathcal{L}_R(r_{\phi}, \mathcal{D}) = -\mathbb{E}_{(x, y_w, y_l)\sim \mathcal{D}}\bigl[\log \sigma(r_{\phi}(x, y_w)- r_{\phi}(x, y_l))\bigr]
\end{equation}

where $\sigma$ denotes the logistic sigmoid. In the setting of language models, $r_\phi(x, y)$ typically leverages weights from the SFT model $\pi^\text{SFT}(y \mid x)$, with a linear head appended atop the last transformer layer, outputting a scalar reward prediction~\cite{ziegler2020finetuning}. To mitigate reward variance, existing methods apply normalization so that $\mathbb{E}_{x,y\sim \mathcal{D}}\left[r_\phi(x, y)\right] = 0$ for each $x$.

\textbf{RL Fine-Tuning Phase}: When entering the RL phase, this trained reward model is employed to deliver evaluative signals for the language model. Consistent with earlier research~\citep{jaques2017sequence, jaques2020human}, the objective can be formalized as

\begin{equation}\label{eq:RL}
\max_{\pi_{\theta}}  \mathbb{E}_{x\sim \mathcal{D}, y\sim \pi_{\theta}(y \mid x)}\bigl[r_{\phi}(x, y)\bigr] - \beta\mathbb{D}_{\textrm{KL}}\bigl[\pi_{\theta}(y\mid x)\mid \mid \pi_\text{ref}(y\mid x)\bigr],
\end{equation}

with $\beta$ acting as a regularization parameter that governs the degree of divergence from a reference policy $\pi_\text{ref}$, which is commonly set as the pre-trained SFT model $\pi^\text{SFT}$. 
Typically, the parameters of the policy $\pi_\theta$ are initialized from $\pi^\text{SFT}$ as well. This KL constraint serves to keep the fine-tuned policy close to the region of input space where the reward model predictions remain reliable and ensures diversity in generation, countering excessive convergence to narrow high-reward responses. Because generating discrete sequences is non-differentiable, reinforcement learning techniques are usually required for optimization. The standard technique~\cite{ziegler2020finetuning, stiennon2022learning, bai2022training, ouyang2022training} introduces the composite reward ${r(x, y) = r_{\phi}(x, y) -\beta (\log \pi_{\theta}(y\mid x) - \log \pi_\text{ref}(y\mid x))}$, and applies PPO~\cite{schulman2017proximal} to maximize it.

\section{Direct Preference Optimization}\label{sec:DPO}

Given the practical difficulties associated with reinforcement learning in large-scale tasks such as language model fine-tuning, we seek an alternative approach for optimizing policies based directly on preference data. Departing from the RLHF paradigm—which separates reward model training and subsequent RL-based policy optimization—we instead capitalize on a reward parameterization that admits a closed-form for the optimal policy, obviating the need for iterative RL. As we elaborate below, the central idea is to use an analytic relationship mapping rewards to policies, permitting reformulation of the reward loss as a policy loss. Through this variable substitution, we sidestep explicit fitting of an isolated reward network, yet remain compatible with popular preference modeling frameworks like Bradley-Terry. Thus, the policy network is tasked with jointly modeling both the generation policy and an implicit reward signal.

\textbf{Deriving the DPO objective.} Consider the standard reinforcement learning objective given by Eq.~\ref{eq:RL}, employing a general reward function $r$. Drawing on earlier research~\citep{peters2007reinforcement, peng2019advantage, korbak2022reinforcement, go2023aligning}, the KL-regularized reward maximization problem in Eq.~\ref{eq:RL} has the following optimal policy form:

\begin{equation}\label{eq:op_policy}
    \pi_r(y\mid x) = \frac{1}{Z(x)}\pi_\text{ref}(y\mid x)\exp\left(\frac{1}{\beta}r(x, y)\right),
\end{equation}

where $Z(x) =\sum_{y}\pi_\text{ref}(y\mid x)\exp\left(\frac{1}{\beta}r(x, y)\right)$ serves as the normalizing constant or partition function. For a comprehensive derivation, refer to Appendix \ref{app:derivation1}. Despite having access to the MLE estimator $r_{\phi}$ of the true reward $r^*$, evaluating $Z(x)$ is computationally costly \citep{korbak2022reinforcement, go2023aligning}, impeding practical usage of this formulation. Nevertheless, one can reformulate Eq.~\ref{eq:op_policy} to write the reward as a function of the optimal policy $\pi_r$, the reference policy $\pi_\text{ref}$, and the unobserved normalization factor $Z(\cdot)$. This is done by taking the logarithm of both sides of Eq.~\ref{eq:op_policy} and rearranging, yielding:

\begin{equation}\label{eq:main_eq}
    r(x,y) =\beta \log \frac{\pi_r(y\mid x)}{\pi_\text{ref}(y\mid x)} + \beta \log Z(x).
\end{equation}

This reparametrization can similarly be applied to the true reward $r^*$ and its optimal policy $\pi^*$. Importantly, the Bradley-Terry model is only a function of the reward difference between two outputs, specifically, ${p^*(y_1 \succ y_2 \mid x) = \sigma(r^*(x, y_1) - r^*(x, y_2))}$. Inserting the above reformulation from Eq.~\ref{eq:main_eq} for $r^*(x, y)$ into the preference probability definition of Eq.~\ref{eq:bradley-terry} causes the partition function terms to cancel, letting us rewrite the preference likelihood purely in terms of $\pi^*$ and $\pi_\text{ref}$. Therefore, under the Bradley-Terry assumption, the optimal RLHF policy $\pi^*$ must satisfy

\begin{equation}\label{eq:objective}
    p^*(y_1\succ y_2 \mid x)=\frac{1}{1 + \exp\left(\beta \log \frac{\pi^*(y_2\mid x)}{\pi_\text{ref}(y_2\mid x)} - \beta \log \frac{\pi^*(y_1\mid x)}{\pi_\text{ref}(y_1\mid x)}\right)}
\end{equation}

A full derivation can be found in Appendix~\ref{app:derivation2}. Although Eq.~\ref{eq:objective} utilizes the Bradley-Terry structure, parallel reasoning extends to broader models such as the Plackett-Luce family~\citep{plackett1975analysis, luce2012individual}, detailed in Appendix~\ref{app:plackett_luce_models}.

Given this new expression for the human preference likelihood framed by the optimal policy, rather than the reward function itself, we can derive a maximum likelihood estimation objective directly over a parametrized policy $\pi_\theta$. Mirroring the procedure used in reward modeling (cf. Eq.~\ref{eq:reward_model}), the corresponding policy-based objective is formulated as:

\begin{equation}\label{eq:optimum_model}
    \mathcal{L}_\text{DPO}(\pi_{\theta}; \pi_\text{ref}) = -\mathbb{E}_{(x, y_w, y_l)\sim \mathcal{D}}\left[\log \sigma \left(\beta \log \frac{\pi_{\theta}(y_w\mid x)}{\pi_\text{ref}(y_w\mid x)} - \beta \log \frac{\pi_{\theta}(y_l\mid

In this manner, we model an implicit reward through an alternative form of parameterization, for which the resulting optimal policy is directly given by $\pi_\theta$. Additionally, as this approach corresponds to optimizing a reparametrized Bradley-Terry model, it possesses well-understood theoretical guarantees, including consistency under appropriate conditions on the distribution of preference data \cite{bong2022generalized}. We provide a deeper examination of these theoretical aspects of DPO and how they compare with related methodologies in Section~\ref{sec:theory}.

\textbf{Interpretation of the DPO update.} To gain mechanistic insight into DPO, it is instructive to inspect the gradient of its loss function, $\mathcal{L}_\text{DPO}$. Specifically, the gradient with respect to $\theta$ can be expressed as follows:

\begin{multline*}\label{eq:gradient}
    \nabla_\theta \mathcal{L}_\text{DPO}(\pi_\theta;\pi_\text{ref}) = \\ -\beta\mathbb{E}_{(x, y_w, y_l) \sim \mathcal{D}} \bigg[\underbrace{\sigma(\hat{r}_\theta(x, y_l) - \hat{r}_\theta (x, y_w))}_\text{greater emphasis when reward assignment is inaccurate}\bigg[\underbrace{\nabla_\theta\log \pi(y_w \mid x)}_\text{promote $y_w$} - \underbrace{\nabla_\theta\log\pi(y_l \mid x)}_\text{penalize $y_l$}\bigg]\bigg],
\end{multline*}

where $\hat{r}_\theta(x, y) = \beta \log \frac{\pi_\theta(y \mid x)}{\pi_\text{ref}(y \mid x)}$ denotes the implicitly specified reward arising from the language model $\pi_\theta$ and the reference model $\pi_\text{ref}$ (see also Section~\ref{sec:theory}). The core idea is that the gradient of $\mathcal{L}_\text{DPO}$ acts to increase the probability assigned to the favored completions $y_w$ while decreasing the probability of the unfavored ones $y_l$. Notably, each sample’s contribution is modulated by how much the implicit reward model $\hat{r}_\theta$ overestimates the less-preferred output, adjusted by the scaling factor $\beta$. This weighting effectively measures the severity of misranking imposed by the current model, in conjunction with the influence of the KL constraint. Empirical results demonstrate the necessity of such weighting: omitting the weighting coefficient, as in a simple variant of this approach, can lead to model collapse or poor generations (refer to Appendix Table~\ref{tab:unlikelihood_generations}).

\textbf{Overview of DPO.} The standard DPO workflow involves the following steps: (1) For each prompt $x$, generate candidate completions $y_1, y_2 \sim \pi_\text{ref}(\cdot \mid x)$, annotate these with human preference judgments to form a dataset of preference pairs $\mathcal{D} = \{x^{(i)}, y_w^{(i)}, y_l^{(i)}\}_{i=1}^N$, and (2) update the language model $\pi_\theta$ by minimizing the DPO objective $\mathcal{L}_\text{DPO}$, conditioned on $\pi_\text{ref}$, $\mathcal{D}$, and the chosen value of $\beta$.

In real-world scenarios, there is often a preference to leverage existing public datasets containing preference data instead of creating new samples and collecting human preferences from scratch. Because these preference datasets are typically generated using $\pi^\text{SFT}$, it is common practice to set $\pi_\text{ref} = \pi^\text{SFT}$ if it is available. On the other hand, if $\pi^\text{SFT}$ is not present, $\pi_\text{ref}$ is obtained by performing maximum likelihood estimation over the preferred completions ${(x, y_w)}$, that is, ${\pi_\text{ref} = \argmax_{\pi}\mathbb{E}_{x, y_w \sim \mathcal{D}}\left[\log \pi(y_w \mid x)\right]}$. This approach serves to reduce the mismatch between the unobserved true reference distribution and the $\pi_\text{ref}$ that DPO actually employs. Readers interested in implementation choices and hyperparameter selection can find additional details in Appendix~\ref{app:implementation}.

\section{Theoretical Analysis of DPO}
In this section, we offer a deeper theoretical interpretation of the DPO approach, present formal justification, and discuss how DPO overcomes some limitations found in actor-critic approaches to RLHF (notably PPO~\cite{schulman2017proximal}).

\label{sec:theory}

\subsection{Your Language Model Is Secretly a Reward Model}
DPO manages to avoid both explicit reward modeling and reinforcement learning, instead leveraging a single maximum likelihood-based objective for learning the policy. Importantly, the optimization objective given in Eq. \ref{eq:main_eq} is equivalent to employing a Bradley-Terry model, where rewards are parameterized as $r^*(x, y) = \beta \log\frac{\pi^*_\theta(y \mid x)}{\pi_\text{ref}(y \mid x)}$. The optimization of the parameterized model $\pi_{\theta}$ thus directly parallels the process of reward model training described in Eq. \ref{eq:reward_model} when expressed through an appropriate variable substitution. In the following, we construct the theoretical foundations underlying this reparameterization, demonstrate that it retains the expressive power of possible reward models, and establish that it is sufficient for recovering the optimal policy. We start by introducing an equivalence relation between reward functions.

\begin{definition}
Two reward functions $r(x, y)$ and $r'(x, y)$ are said to be equivalent if and only if ${r(x, y)-r'(x, y) = f(x)}$ for some function $f$.
\end{definition}

This definition straightforwardly gives rise to an equivalence relation, thereby partitioning the set of all reward functions into distinct equivalence classes. We now present two key lemmas:

\begin{lemma}\label{lemma:same_prefrence} Within the Plackett-Luce model, including the Bradley-Terry model as a special case, any two reward functions belonging to the same equivalence class generate identical preference distributions.
\end{lemma}

\begin{lemma}\label{lemma:same_policy}
    For the constrained RL problem, any two reward functions in the same equivalence class will produce the same optimal policy.
\end{lemma}

Since the proofs are direct, we postpone them to Appendix \ref{app:lemma1}. The initial lemma captures a well-documented under-identifiability issue associated with the Plackett-Luce family of models \cite{plackett1975analysis}. This lack of identifiability necessitates imposing additional constraints when seeking meaningful guarantees for the MLE estimators of Eq. \ref{eq:reward_model} \cite{bong2022generalized}. The second lemma clarifies that, since the optimal policy is invariant across equivalent reward functions, the precise choice of reward function within an equivalence class is inconsequential for optimal control; thus, it suffices to identify any member from the optimal class. We establish the following theorem in Appendix~\ref{app:thm1}:

\begin{theorem}\label{thm:main}
    Given mild technical assumptions, every equivalence class of reward functions consistent with Plackett-Luce (and in particular, Bradley-Terry) models admits the representation ${r(x, y) = \beta \log \frac{\pi(y\mid x)}{\pi_\text{ref}(y\mid x)}}$ for some candidate model $\pi(y\mid x)$ and a chosen reference distribution $\pi_\text{ref}(y \mid x)$.
\end{theorem}

\begin{sproof}
    Let $r(x, y)$ be any reward function, with associated optimal distribution $\pi_r(y \mid x)$ as prescribed by Eq. \ref{eq:op_policy}. We demonstrate that some member of the equivalence class of $r$ can always be expressed in the form specified above. To do so, we introduce the following projection:
\begin{equation}
    f(r; \pi_\text{ref}, \beta)(x, y) = r(x, y) - \beta\log\sum_{y}\pi_\text{ref}(y\mid x)\exp\left(\frac{1}{\beta}r(x, y)\right)
\end{equation}
The transformation $f$ applies a normalization by subtracting the logarithm of the partition function with respect to $\pi_r$. Since this normalization is only dependent on $x$, $f(r; \pi_\text{ref}, \beta)(x, y)$ remains in the same equivalence class as $r(x, y)$. Upon substituting the reward as given on the right-hand side of Eq.~\ref{eq:main_eq} (which is universally valid), it follows that $f(r; \pi_\text{ref}, \beta)(x, y) = \beta \log \frac{\pi_r(y\mid x)}{\pi_\text{ref}(y\mid x)}$. Therefore, this mapping yields an element in the same equivalence class with the required structure, establishing that the proposed reparameterization captures all relevant reward functions without loss of generality.
\end{sproof}

An alternative interpretation of Theorem~\ref{thm:main} is that it characterizes the specific representative reward function, within each equivalence class, chosen by the DPO reparameterization—namely, the reward function that fulfills:

\begin{equation}\label{eq:lag_p}
     \sum_{y}\underbrace{\pi_\text{ref}(y\mid x)\exp\left(\frac{1}{\beta}r(x, y)\right)}_{=\pi(y\mid x)\text{, using Thm.~\ref{thm:main} reparam.}} = 1,
\end{equation}

which ensures $\pi(y\mid x)$ forms a legitimate probability distribution (all probabilities are nonnegative and sum to unity). Notably, as seen from Eq.~\ref{eq:op_policy}, this condition corresponds to the partition function associated with the optimal policy determined by the reward function $r(x, y)$. The primary innovation in the DPO approach is recognizing that one can impose suitable constraints on the otherwise underspecified Plackett-Luce family of preference models (in particular, the Bradley-Terry model), thereby maintaining the expressivity of the reward model class while making the optimal policy from Eq.~\ref{eq:op_policy} analytically solvable for arbitrary prompts $x$.

\subsection{Instability of Actor-Critic Algorithms} Our framework also provides insight into sources of instability in standard actor-critic algorithms widely employed in RLHF pipelines, including PPO. Adhering to the standard RLHF workflow, we analyze the reinforcement learning fine-tuning stage described in Section \ref{section:prelims}. The approach is connected to the control-as-inference formalism \cite{levine2018reinforcement} when applied to the constrained RL setup as formulated in \ref{eq:RL}. Consider a parameterized policy $\pi_{\theta}(y\mid x)$; the objective becomes minimizing the Kullback-Leibler divergence $\mathbb{D}_{\text{KL}}[\pi_{\theta}(y|x) \mid \mid \pi^*(y\mid x)]$, where $\pi^*$ is the optimal policy defined by Eq.~\ref{eq:optimum_model} and shaped by the reward function $r_{\phi}(y, x)$. Elementary manipulations yield the following objective for optimization:

\begin{equation}\label{eq:AC}
    \max_{\pi_{\theta}}\mathbb{E}_{\pi_{\theta}(y\mid x)}\bigg[\underbrace{r_{\phi}(x, y) -\beta\log\sum_{y}\pi_\text{ref}(y\mid x)\exp\left(\frac{1}{\beta}r_{\phi}(x, y)\right)}_{f(r_{\phi}, \pi_\text{ref}, \beta)} - \underbrace{\beta\log\frac{\pi_{\theta}(y\mid x)}{\pi_\text{ref}(y\mid x)}}_{\text{KL}}\bigg]
\end{equation}

This objective mirrors that used in earlier studies \citep{ziegler2020finetuning, stiennon2022learning, bai2022training, ouyang2022training}, where optimization proceeds via a DPO-equivalent reward within the class of rewards parameterized by $r_{\phi}$. Within this framework, the normalization term present in $f(r_{\phi}, \pi_\text{ref}, \beta)$ can be interpreted as the reference policy $\pi_\text{ref}$’s soft value function. Although this normalization does not alter the location of the optimal policy, omitting it may result in high-variance policy gradient estimates, thereby destabilizing training. To address this, a learned value function can be incorporated to approximate the normalization, though optimizing such an auxiliary function may also be challenging. Previous works have addressed normalization by anchoring rewards to a human completion, which amounts to using a single-sample Monte Carlo estimate for the normalization. By contrast, under the DPO reparameterization, the constructed reward function circumvents the need for such baselines.

\section{Experiments}
Here, we empirically assess the capability of DPO to train policies using preference data directly. We begin by investigating DPO in a carefully controlled text-generation scenario, asking: how effectively does DPO balance the objectives of reward maximization and minimizing KL-divergence relative to the reference policy, particularly when compared with prevalent preference optimization approaches like PPO? We then extend our study to consider DPO’s behavior on larger models and more challenging RLHF tasks, encompassing both summarization and dialogue. Our results indicate that, with minimal hyperparameter adjustment, DPO typically achieves performance on par with or superior to strong baselines, such as PPO-based RLHF, and surpasses approaches selecting the top-performing trajectory among $N$ samples with respect to a learned reward. We first outline the experimental methodology, with further specifics available in Appendix~\ref{app:exp_details}.

\textbf{Tasks.} Our study investigates three distinct open-ended text generation tasks. In each scenario, the employed algorithms derive a policy using a preference dataset, denoted as $\mathcal{D}=\bigl\{x^{(i)}, y_w^{(i)}, y_l^{(i)}\bigr\}_{i=1}^N$. In the task of \textbf{controlled sentiment generation}, $x$ represents the initial portion of a movie review sourced from the IMDb dataset \cite{maas-EtAl:2011:ACL-HLT2011}, and the policy's objective is to generate a continuation $y$ that exhibits positive sentiment. To systematically assess performance, we construct pairs of generated texts for comparison using a sentiment classifier pre-trained for this purpose, such that $p(\text{positive}\mid x,y_w)>p(\text{positive}\mid x,y_l)$. The SFT baseline involves fine-tuning the GPT-2-large model on the training set of IMDb reviews until performance stabilizes (see App~\ref{app:sentiment_details} for additional details). For \textbf{summarization}, $x$ consists of a Reddit forum post, and the model must create a summary $y$ highlighting its main ideas. Following established methods, we utilize the Reddit TL;DR summarization dataset \citep{volske-etal-2017-tl}, accompanied by human preference judgments as compiled by \citeauthor{stiennon2022learning}. Our SFT baseline is produced by further training a model on human-generated Reddit summaries,\footnote{\url{https://huggingface.co/CarperAI/openai_summarize_tldr_sft}} employing the TRLX \citep{leandro_von_werra_2023_7790115} framework for RLHF. Notably, the preference data in this case were collected on outputs from an alternate SFT model trained similarly. Lastly, for \textbf{single-turn dialogue}, $x$ constitutes a user's input, varying from technical questions to personal advice inquiries. The goal is for the policy to generate a response $y$ that is both useful and engaging; to this end, we leverage the Anthropic Helpful and Harmless dialogue dataset \citep{bai2022training}, which contains 170k exchanges between human participants and an automated assistant. Each dialogue ends with two alternative assistant responses produced by a large language model, accompanied by a label indicating which one is preferred by a human annotator. As no pre-trained SFT model is provided in this context, we develop an SFT baseline by fine-tuning an available language model exclusively on the preferred response completions.

\textbf{Evaluation.} We employ two complementary evaluation strategies in our experiments. To specifically assess each algorithm’s performance on the constrained reward maximization objective within the controlled sentiment generation scenario, we examine the trade-off frontier between achieved reward and KL-divergence from the reference policy; this analysis is feasible since the actual reward function (i.e., a sentiment classifier) is known. Conversely, because the true reward function is not accessible in real-world conditions, we rely on \textit{win rate} comparisons against a baseline policy, leveraging GPT-4 as a stand-in for human evaluators to judge summary quality in summarization and helpfulness in single-turn dialogue tasks. For summarization, the baseline is the reference summary provided in the test set; for dialogue, it is the test set’s preferred response. Although recent literature suggests that large language models can serve as superior automatic evaluators compared to existing metrics \citep{Chen2023ExploringTU}, we also present a human-subject study to substantiate our selection of GPT-4 for evaluation, as described in Sec.~\ref{sec:human-judgments}. Our findings demonstrate a strong correlation between GPT-4 and human assessments, with levels of human-GPT-4 agreement that are comparable to, or surpass, inter-annotator reliability.

\textbf{Methods.} Alongside DPO, we benchmark a variety of established methods for aligning language model outputs with human preferences. As baseline prompting strategies, we examine zero-shot prompting with \textbf{GPT-J} \citep{gpt-j} for summarization and 2-shot prompting with \textbf{Pythia-2.8B} \citep{biderman2023pythia} for dialogue. Additionally, we include the \textbf{SFT} model and introduce \textbf{Preferred-FT}, which undergoes supervised fine-tuning on the chosen outputs $y_w$, sourced from either SFT (for sentiment control and summarization) or a general language model (for dialogue). Another comparable supervised approach is \textbf{Unlikelihood}~\citep{welleck2019neural}, in which the model is trained to enhance the likelihood of $y_w$ while suppressing the likelihood of $y_l$, modulated by an optional weighting factor $\alpha\in[0,1]$ for the unlikelihood loss. Furthermore, we incorporate \textbf{PPO} \citep{schulman2017proximal}, where the reward model is learned from preference annotations, and \textbf{PPO-GT}, which acts as an oracle policy by training directly with the ground-truth reward (available in sentiment-controlled experiments). For sentiment tasks, our setup includes both a standard implementation of PPO-GT \cite{leandro_von_werra_2023_7790115} and an enhanced version that applies reward normalization and additional hyperparameter tuning for improved results (these improvements are similarly applied to PPO with learned rewards). Lastly, we evaluate a \textbf{Best of $N$} strategy, where $N$ samples are drawn from the SFT model (or Preferred-FT in the case of dialogue), and the candidate receiving the highest score from a learned reward model is chosen. While this approach is capable of achieving strong performance by dissociating the influence of reward model accuracy from PPO training, it becomes infeasible for larger $N$ due to the requirement of generating $N$ completions per test instance.

\subsection{How well can DPO optimize the RLHF objective?}

The reward maximization objective in standard RLHF algorithms includes a KL constraint to balance reward exploitation with limiting the extent to which the policy diverges from a reference policy. As such, a fair evaluation of different methods necessitates considering both the reward obtained and the magnitude of the KL divergence; obtaining modest gains in reward at the expense of much larger KL increases is not necessarily favorable. In Figure~\ref{fig:frontier-tldr-main}, we display the tradeoff between reward and KL divergence for several approaches in the sentiment task. For each algorithm, we conduct several training runs, varying the conservativeness parameter in each run (target KL $\in\{3,6,9,12\}$ for PPO; $\beta \in \{0.05,0.1,1,5\}$ for DPO; $\alpha\in\{0.05,0.1,0.5,1\}$ for unlikelihood; as well as using different random seeds for preferred-FT), comprising a total of 22 runs. After every 100 steps up to convergence, each trained policy is assessed on a test set, computing the average true reward and the mean sequence-level KL\footnote{Namely, the sum over all timestep-wise KL divergences.} relative to the reference policy, i.e., $\text{KL}\left(\pi\mid \mid \pi_\text{ref}\right)$. The results reveal that DPO establishes the most favorable reward-KL frontier, attaining the highest reward for comparatively small KL divergence. This finding is significant for several reasons. Most notably, although DPO and PPO target the same objective, DPO outperforms PPO with respect to reward efficiency per unit of KL, offering a superior reward/KL tradeoff. Additionally, DPO surpasses PPO even in settings where PPO is given access to ground truth rewards (PPO-GT).

\subsection{Can DPO scale to real preference datasets?}
\label{sec:dpo-real-datasets}

We proceed to assess how DPO performs in fine-tuning for tasks involving summarization and single-turn dialogue. In the case of summarization, conventional automatic metrics like ROUGE have been shown to exhibit weak alignment with human judgments~\citep{stiennon2022learning}. Previous studies indicate that preference-based fine-tuning of language models using PPO can yield summaries that better align with human preferences. To compare methods, we generate outputs on the TL;DR summarization dataset's test split, calculating the mean win rate over the reference summaries. For each method, outputs are sampled across a temperature range from 0.0 to 1.0; the corresponding win rates are illustrated in Figure~\ref{fig:frontier-tldr-main} (right). All models—DPO, PPO, and Preferred-FT—are fine-tuned starting from the same GPT-J SFT checkpoint\footnote{\url{https://huggingface.co/CarperAI/openai_summarize_tldr_sft}}. Our results indicate that DPO achieves a win rate of roughly 61\% at temperature 0.0, outperforming PPO, which attains about 57\% at its own optimal sampling temperature of 0.0. DPO additionally surpasses the best win rate of the $N$-best baseline. We emphasize that the $\beta$ hyperparameter for DPO was not systematically optimized, suggesting these results may be conservative estimates of its full capability. Furthermore, DPO displays substantially greater robustness across different sampling temperatures compared to PPO, whose effectiveness may decrease to baseline GPT-J levels at higher temperatures. In contrast, Preferred-FT yields only marginal improvements over SFT. We further conduct direct human preference comparisons between DPO and PPO in Section~\ref{sec:human-judgments}, finding that human raters preferred samples from DPO at temperature 0.25 over those from PPO at temperature 0 in 58\% of cases.

For single-turn dialogue assessment, we focus on the subset of the test partition from the Anthropic HH dataset \citep{bai2022training} that involves just one exchange between a human and the assistant. The evaluation with GPT-4 leverages the preferred completions as references for determining the win rate among the various methods. Since an established SFT model for this specific task does not exist, our approach begins with a pre-trained Pythia-2.8B model; we first use Preferred-FT to develop a reference model by training on the selected completions so that outputs remain aligned with the model’s distribution. Subsequently, DPO training is performed. Additionally, we compare DPO to both the strongest result from 128 Preferred-FT completions (having observed that the Best of $N$ approach plateaus at $N = 128$; refer to Appendix Figure~\ref{fig:best-of-n}) and a version of Pythia-2.8B base employing a 2-shot prompt, noting that DPO achieves equal or superior results at optimal temperature settings for each method. An RLHF model, trained using PPO on the Anthropic HH dataset \footnote{\url{https://huggingface.co/reciprocate/ppo_hh_pythia-6B}}, is also evaluated, as provided by a reputable repository \footnote{\url{https://github.com/CarperAI/trlx/tree/main/examples/hh}}; however, we do not identify a combination of prompt and sampling temperature that surpasses the original Pythia-2.8B model in performance. Given our findings on TL;DR and recognizing that Best of 128 and PPO both optimize an identical reward, we regard Best of 128 as a loose proxy for PPO-level outcomes. In summary, DPO stands out as the only method that is computationally efficient and delivers performance gains over the dataset’s preferred completions, matching or outperforming the much more resource-intensive Best of 128 benchmark. Furthermore, Figure~\ref{fig:dialogue-main} demonstrates that DPO reaches optimal performance rapidly.

\subsection{Generalization to a new input distribution}

To further assess how PPO and DPO respond to distributional changes, we tested the policies trained in the Reddit TL;DR summarization setup on a distinct dataset, namely news articles from the CNN/DailyMail test set \citep{nallapati-etal-2016-abstractive}, and utilized the optimal sampling temperatures identified for TL;DR (0 and 0.25). Results, shown in Table~\ref{tab:ood}, indicate that GPT-4 win rates—computed relative to ground-truth dataset summaries—were determined using the identical GPT-4 (C) prompt as in Reddit TL;DR evaluations, with ``forum post'' substituted for ``news article''. On this shifted input domain, DPO maintained a marked lead over PPO. These findings suggest that DPO policies possess robust generalization capabilities on par with PPO, even though DPO does not leverage the additional unlabeled Reddit TL;DR prompts available to PPO.

\subsection{Validating GPT-4 judgments with human judgments}
\label{sec:human-judgments}
A human study was conducted to assess the consistency of GPT-4’s evaluation, drawing from outcomes of the TL;DR summarization experiment and utilizing two distinct GPT-4 prompts. The \textbf{GPT-4 (S)} (simple) prompt requests a determination of which summary best captures the core information of the post. By contrast, the \textbf{GPT-4 (C)} (concise) prompt additionally inquires about conciseness, a feature we examined given that GPT-4, when using the \textbf{GPT-4 (S)} prompt, tends to prefer more verbose and repetitive summaries than human evaluators. Complete prompt texts can be found in Appendix~\ref{app:prompts}. Three comparative analyses were performed, spanning the highest performing method (DPO, temp. 0.25), the lowest (PPO, temp. 1.0), and an intermediate method (SFT, temp. 0.25) to capture a spectrum of sample qualities; each was compared against greedily-sampled PPO at its optimal temperature. Across both prompts, GPT-4’s evaluations aligned with human preferences to a degree similar to inter-human agreement, indicating that GPT-4 serves as a reasonably faithful proxy for human assessment (for DPO and PPO-1, we collected multiple independent human judgments due to limited human resources). In summary, the \textbf{GPT-4 (C)} prompt generally yields win rates most reflective of human opinion, motivating its use for the principal results in Section~\ref{sec:dpo-real-datasets}. Further details regarding the human study methodology, web interface, and participant roster are available in Appendix~\ref{app:human-study}.

\section{Discussion}
Preference-based learning offers a robust and extensible approach for developing advanced, aligned language models. In this work, we have presented DPO, a streamlined methodology that facilitates preference-driven language model training without the need for reinforcement learning. Instead of adapting the preference optimization challenge to a conventional RL framework to leverage standard RL algorithms, DPO establishes a direct connection between language model policies and reward functions. This allows for optimizing language models to meet human preferences through a straightforward cross-entropy objective—eschewing both reinforcement learning and any loss of expressivity. DPO requires little hyperparameter adjustment and attains results that are comparable to or surpass existing RLHF approaches, notably those employing PPO. Consequently, DPO substantially lowers the hurdles associated with preference-based language model training.

\textbf{Limitations \& Future Work.} Several avenues for further investigation are prompted by our findings. One open question concerns the extent to which the DPO approach can generalize beyond its training distribution, especially relative to methods utilizing explicit reward functions. Preliminary evidence indicates that DPO exhibits similar generalization to PPO-derived models, yet systematic evaluation remains outstanding. For instance, it is unclear whether the use of self-labeling strategies with DPO can capitalize on unlabeled input prompts to the same degree. Additionally, understanding how over-optimization of reward might emerge in direct preference optimization, as suggested by the slight reduction in effectiveness depicted in Figure~\ref{fig:dialogue-main}-right, deserves closer examination. While our experiments encompass models up to 6B parameters, scaling DPO to models with substantially larger capacities constitutes a promising area for subsequent research. In terms of assessment, our observations reveal that GPT-4-based win rates can vary depending on the prompt design; further study into optimizing automated evaluation quality is warranted. Finally, the DPO methodology offers potential beyond the realm of language model training from human preferences, with prospective applications in training generative models across other domains.